% Appendix: the frontier experiment. Moved from Section~\ref{sec:qual-frontiers}. Source:
% evidence/x2-x5-run1.md (GitHub Actions run 40921769714, commit 884ce18; pre-registration
% PREREG_X2_X5.md at 828ea27), every number as x2x5_stats.py printed it; QPS from each store's
% build-0 job.

\section{The frontier experiment in detail (X2)}\label{app:frontier}

This appendix gives the design and full results of the frontier experiment, which Section~\ref{sec:qual-frontiers}
summarizes. Its verdicts are as the analysis script printed them (Appendix~\ref{app:ledger}).

The enterprise scenario has one permission layout. To separate the filter's selectivity from its
correlation with the question, X2 adds a corpus whose layouts we control: the first million
passages of MS MARCO~\cite{bajaj2016ms} from BEIR, embedded with bge-small-en-v1.5 (384 dimensions,
so that ten stores fit beside the harness on a 16\,GB runner). Every passage belongs to one team of
a hierarchy of 100 divisions, 10 departments each and 10 teams each, assigned either at random
(\emph{uncorrelated}) or by hierarchical $k$-means over the embeddings (\emph{clustered}, so groups
are topical), five seeds of each. A subject receives divisions, departments and teams until it may
read a target share of the corpus, from 0.06\% to 100\%, and its filter is either one disjunction over
its groups or a tenant clause conjoined with a disjunction over a grant twice as large, a shape in
which each clause alone passes far more than both. Questions come from a pool of 20,000 MS MARCO
questions: four drawn at random for a subject of an uncorrelated layout; for a subject of a clustered
layout, four from those whose 100 nearest passages are at least half readable (\emph{positively
correlated}) and four from those with nothing readable among their 1,000 nearest (\emph{negatively
correlated}), the best by the criterion when fewer than four qualify (684 of the 6,200 pairs). Ten subjects per seed and cell make 1,050
subjects. For every question and subject we compute, from exact scores, the global selectivity,
the local selectivity among the 100 and 1,000 nearest passages, the rank of the first readable
passage, ACORN's query correlation~\cite{patel2024acorn}, the global-local selectivity of Amanbayev
et al.~\cite{amanbayev2026filtered} and the over-fetch estimate of HCBGen~\cite{lim2026revisiting}.

Ten approximate stores run against FAISS over a flat index: pgvector, Qdrant, Elasticsearch,
OpenSearch, Weaviate, Redis, Milvus, MongoDB, LanceDB and Chroma, each at its default, at five levels
of its search-effort parameter, with its store-specific options (pgvector's raised scan limits and
strict order, Qdrant's ACORN and per-tenant graphs, Weaviate's filter strategy and flat-search cutoff,
Redis's hybrid policies, LanceDB's refinement and post-filtering, and on the enterprise corpus
Elasticsearch's \texttt{int8\_hnsw} index beside its default \texttt{bbq\_hnsw}) and with its own
exact search where it has one (not Milvus or Chroma), in three fresh builds per store; exact searches,
which do not use the index, run on the first build, and FAISS on one. On MS MARCO the per-tenant
graphs are queried by the subjects with a tenant clause. Each pgvector search is planned with its own
values, so the plan-cache effect of Section~\ref{sec:qual-pgvector} cannot arise, and the share of
searches whose plan uses the HNSW index is recorded. Throughput is measured end to end through the
layer with one client thread and compared only within a job; recall is compared across jobs. Each
setting has a share of its job's time; one that runs out keeps a prefix of a fixed order that cycles
through cells and subjects. The same harness runs the enterprise scenario, whose as-run settings must
reproduce the audit's recall within 0.03 for seven of nine stores. We pre-registered four hypotheses: that negatively correlated questions lose at least 0.05
of recall@10 against uncorrelated ones at defaults in at least six stores (H2.1); that local
selectivity predicts recall better than global selectivity (H2.2); that raising search effort
narrows but does not close that gap (H2.3); and that intermediate selectivities are harder than very
small ones, which stores answer by brute force (H2.4).

\paragraph{What ran.}
X2 ran on GitHub Actions hosted runners (four CPUs, 16\,GB), one store per job: PostgreSQL 17.11
with pgvector 0.8.7 (the audit ran PostgreSQL 16.15), Qdrant 1.19.2, Elasticsearch 9.5.5, OpenSearch 3.9.0, Weaviate 1.39.12, Redis
8.10.2, Milvus (image 3.0.2), MongoDB Atlas Local 8.0.30, LanceDB 0.40.0 and Chroma 1.5.9. The
validity checks held. Every exact setting, and FAISS, returned the exact answer to every question it
ran (recall@10 1.000); no hit was unreadable and no canary appeared; and every subject granted 0.5\%
or more could read within 25\% of its target share. Eight of the nine stores of the audit reproduced
the recall@10 of the first execution of its seed-7 run, the values recorded beside the X2 verdicts in \path{evidence/x2-x5-run1.md}
(not the four-run means of Table~\ref{tab:qual-x0}), within 0.03 at its
settings, pgvector among
them at 0.730 against 0.732, with every search planned with its own values; against the execution the artifact holds
(Appendix~\ref{sec:e2e-setup}) seven of the nine are within 0.03, Qdrant's 0.576 against 0.530 being the other exception, which
still meets the pre-registered rule of at least seven. Chroma reached 0.990
against 0.883 (0.878 in the artifact's execution): at about 1\,s a retrieval its share of the job's time covered only 381 to 407 of the
1,468 questions, the first questions of each person, which for most people ask about a document they
may read, the kind it answered best in the audit (overlap@5 of 0.989, against 0.82 to 0.84 for the
other kinds), a reading consistent with the gap though the run did not print recall by kind. On
MS MARCO Chroma was not measured, a deviation from the plan. Each of its three builds loaded the million
passages in 22 to 28 minutes, but a filtered retrieval then took 137 to 221\,s, so the twenty
warm-up retrievals of each setting used up the setting's time and none reached a timed pass; under
the pre-registered procedure a rerun cannot finish within a job.

\begin{table}[tbp]
\centering
\footnotesize
\setlength{\tabcolsep}{3.5pt}
\begin{tabular}{@{}lrrrlrrlrrrr@{}}
\toprule
 & \multicolumn{3}{c}{Default} & \multicolumn{3}{c}{Most search effort} & \multicolumn{3}{c}{Best other option} & Exact & Ent. \\
\cmidrule(lr){2-4}\cmidrule(lr){5-7}\cmidrule(lr){8-10}
Store & R@10 & 100\% & QPS & Setting & R@10 & QPS & Option & R@10 & QPS & QPS & R@10 \\
\midrule
pgvector & 0.575 & 0.916 & 75.3 & \texttt{ef\_search} 1,000 & 0.927 & 14.8 & both scan limits & 0.949 & 12.9 & 0.26 & 0.112 \\
Qdrant & 0.799 & 0.978 & 205.3 & \texttt{hnsw\_ef} 1,024 & 0.898 & 96.7 & ACORN 1.0 & 0.997 & 28.5 & 22.6 & 0.641 \\
Elasticsearch & 0.806 & 0.870 & 185.1 & candidates 10,000 & 0.878 & 107.2 & & & & 52.1 & 0.815 \\
OpenSearch & 0.998 & 0.971 & 105.3 & \texttt{ef\_search} 512 & 0.999 & 91.5 & & & & 15.8 & 0.948 \\
Weaviate & 0.968 & 0.966 & 136.1 & \texttt{ef} 1,024 & 0.994 & 77.1 & sweeping & 0.995 & 104.3 & 24.8 & 0.923 \\
Redis & 0.971 & 0.793 & 54.1 & \texttt{EF\_RUNTIME} 1,000 & 0.999 & 49.9 & & & & 17.0 & 0.987 \\
Milvus & 0.972 & 0.936 & 68.6 & \texttt{ef} 512 & 0.999 & 98.1 & & & & & 0.864 \\
MongoDB & 0.996 & 0.960 & 24.1 & candidates 10,000 & 1.000 & 10.3 & & & & 18.7 & 0.953 \\
LanceDB & 0.437 & 0.517 & 67.9 & \texttt{nprobes} 500 & 0.450 & 24.1 & refine 30 & 0.908 & 37.9 & 4.6 & 0.439 \\
Chroma & \multicolumn{10}{c}{not measured: 137 to 221\,s per retrieval} & 0.983 \\
FAISS (flat) & & & & & & & & & & 0.91 & 1.000 \\
\bottomrule
\end{tabular}
\caption{X2 on MS MARCO (1,000,000 passages; 6,200 question--subject pairs over every layout,
correlation regime and selectivity). R@10: recall@10, the mean over three builds (exact settings:
one). 100\%: recall@10 at the default on the uncorrelated subjects granted the whole corpus, whose
filter lists every division and so passes every passage (200 pairs per build, the mean over three
builds), the nearest the data come to an unfiltered baseline. QPS: one client thread through the
layer, the median of three passes in the store's build-0
job, so comparable only within a row. Most search effort: the top of the store's five-level sweep.
pgvector's default has iterative scans off; ``both scan limits'' raises \texttt{hnsw.max\_scan\_tuples}
to 1,000,000 and \texttt{hnsw.scan\_mem\_multiplier} to 16 at \texttt{ef\_search} 100; its exact
search answered 998 pairs in its time. Qdrant's ACORN option has \texttt{max\_selectivity} 1.0 and
\texttt{hnsw\_ef} 512. Ent.: recall@10 at the default on the enterprise corpus (1,468 questions;
Elasticsearch's \texttt{bbq\_hnsw}, 0.888 with \texttt{int8\_hnsw}; Chroma on 407).}
\label{tab:qual-frontiers}
\end{table}

\paragraph{Defaults.}
At the settings a user gets without tuning, recall@10 on MS MARCO ranged from 0.437 and 0.575 to
0.998 (Table~\ref{tab:qual-frontiers}). The two lowest are documented defaults. pgvector without
iterative scans filters the 40 candidates its index returns and stops: 32.4\% of the retrievals on MS
MARCO and 61.6\% on the enterprise corpus came back empty (recall@10 0.575 and 0.112), against 0.3\%
and 9.7\% with the iterative scans the layer configures (recall@10 0.897 and 0.661 at
\texttt{ef\_search} 100). LanceDB's default IVF-PQ index ranks by quantized distances: raising
\texttt{nprobes} from 10 to 500 moved recall@10 only from 0.424 to 0.450 at a third of the throughput,
while re-ranking 30 times $k$ candidates by their full vectors reached 0.908. Elsewhere effort bought
recall at a price in throughput, and the stores' own options bought more: Qdrant reached 0.898 at the
top of its sweep but 0.997 with ACORN, and Weaviate's sweeping filter strategy reached 0.995 against
0.968 for its default ACORN strategy~\cite{weaviate2026vectorindex}, at 104.3 against 136.1 queries
per second. Index builds varied little except Qdrant's: its recall@10 at the default ranged from
0.759 to 0.878 over three builds of the same data, as in X1, wider than the intervals below, which resample subjects and are conditional on the builds.

To separate the filter's cost from the index's own error, Table~\ref{tab:qual-frontiers} also gives
recall@10 at 100\% selectivity, where the filter passes everything. At the defaults that baseline
ran from 0.517 for LanceDB, the floor of its quantized index, and 0.793 for Redis, at an
\texttt{EF\_RUNTIME} of 10, to 0.978 for Qdrant. Four stores lost recall against it over the whole
mix of filters: pgvector (0.575 against 0.916), Qdrant (0.799 against 0.978), Elasticsearch (0.806
against 0.870) and LanceDB (0.437 against 0.517). The other five scored no lower than at 100\%,
consistent with many of the mix's filters being small enough for the store to answer by brute force;
what they lost at 100\% is the index's error, not the filter's. The enterprise corpus has no such
cell, because no person may read all of it.

Read at a common throughput, the same sweeps need a caveat: each store's queries per second come from
its own job on its own runner, so they are not calibrated across stores. Among the settings that
served at least 50 queries per second in their build-0 job, exact search included, the best recall@10
of that build was 1.000 for Elasticsearch (its exact search, at 52; 0.878 for its approximate search
at 107), 0.999 for OpenSearch and Milvus,
0.997 for Redis, 0.995 for Weaviate, 0.979 for Qdrant (ACORN at its default gate, at 69), 0.578 for
pgvector (its default without iterative scans; with them it served at most 16) and 0.438 for LanceDB;
MongoDB reached 50 at no setting (at most 28). The threshold of 50 is this reading's and not the pre-registration's. A calibrated comparison would need every store timed
on one machine, or a reference workload on each runner, and interpolation along each sweep.

\begin{table}[tbp]
\centering
\scriptsize\setlength{\tabcolsep}{3pt}
\begin{tabular}{@{}llllll@{}}
\toprule
Store & \multicolumn{1}{c}{$G$ (H2.1)} & \multicolumn{1}{c}{$D$ (H2.2)} & & \multicolumn{1}{c}{$V$ (H2.4)} & \\
\midrule
pgvector & 0.4242 (0.4154, 0.4332) gap & 0.3650 (0.3321, 0.3989) & local & 0.9267 (0.8867, 0.9600) & harder \\
Qdrant & 0.0343 (0.0196, 0.0487) & $-0.2032$ ($-0.2194$, $-0.1876$) & global & 0.5403 (0.5165, 0.5653) & harder \\
Elasticsearch & 0.1371 (0.1242, 0.1500) gap & 0.2439 (0.2116, 0.2754) & local & 0.2338 (0.2032, 0.2650) & harder \\
OpenSearch & $-0.0014$ ($-0.0029$, 0.0001) none & $-0.0601$ ($-0.0793$, $-0.0411$) & n.a. & 0.0005 (0.0000, 0.0015) & \\
Weaviate & 0.0549 (0.0478, 0.0618) gap & $-0.3900$ ($-0.4159$, $-0.3644$) & global & 0.1780 (0.1578, 0.1980) & harder \\
Redis & $-0.0794$ ($-0.0871$, $-0.0719$) & $-0.0757$ ($-0.0870$, $-0.0647$) & global & 0.0000 (0.0000, 0.0000) & not harder \\
Milvus & 0.0136 (0.0080, 0.0194) none & $-0.5307$ ($-0.5469$, $-0.5140$) & global & 0.1037 (0.0925, 0.1150) & harder \\
MongoDB & $-0.0037$ ($-0.0055$, $-0.0019$) none & $-0.0804$ ($-0.1016$, $-0.0593$) & global & 0.0012 (0.0005, 0.0020) & \\
LanceDB & 0.2535 (0.2423, 0.2649) gap & 0.5382 (0.5080, 0.5674) & local & $-0.0098$ ($-0.0332$, 0.0212) & not harder \\
\bottomrule
\end{tabular}
\caption{X2's pre-registered statistics on MS MARCO at the stores' defaults, with 95\% intervals from
a cluster bootstrap over subjects (10,000 resamples). $G$: recall@10 of uncorrelated minus negatively
correlated questions at selectivities 2\% to 30\% (``gap'' needs $G \ge 0.05$, an interval above 0 and
the same sign in at least four of five seeds; ``none'' $|G| < 0.02$ with the interval inside
$\pm 0.03$). $D$: $|\rho_S(\text{recall@10}, \rho_{1000})| - |\rho_S(\text{recall@10}, \sigma_g)|$,
Spearman over the 6,200 pairs (``local'' $D \ge 0.05$ with the interval above 0, ``global'' $D < 0$
with it below; n.a.: fewer than 100 pairs with recall below 1). $V$: recall@10 at 0.06\% minus the
lowest at 0.5\%, 2\% and 10\%, clustered layouts, negatively correlated questions, one disjunction
(``harder'', intermediate filters harder than the smallest, needs $V \ge 0.05$, an interval above 0
and four seeds; ``not harder'' $V \le 0$; the script printed these as ``U-shaped'' and ``monotone'',
but $V$ does not test recovery at larger selectivities). A blank
label is neither outcome. Chroma was not measured.}
\label{tab:qual-x2tests}
\end{table}

\paragraph{Correlation (H2.1, inconclusive).}
Negatively correlated questions lost at least 0.05 of recall@10 in four stores, not the six we had
predicted, and no loss beyond 0.02 in three (Table~\ref{tab:qual-x2tests}): pgvector lost 0.424,
LanceDB 0.254, Elasticsearch 0.137 and Weaviate 0.055, each in all five seeds; OpenSearch, MongoDB and
Milvus stayed within 0.014 at recall@10 of 0.97 to 1.00. The other two show why one number per store
misleads. Redis's recall was \emph{higher} for negatively correlated questions, by 0.079, and the cost
appeared as latency: in the build-0 job those questions at 10\% and 30\% selectivity were answered
exactly at a median of 56.8 and 147.1\,ms, against 0.908 and 0.858 of recall at 2.2 and 2.7\,ms for
uncorrelated ones, consistent with Redis switching from batches of graph neighbors to brute force over
the filter's matches when the batches find too few~\cite{redis2026vectors}. OpenSearch and MongoDB
paid the same way, more mildly (24.6 against 15.3\,ms, and 68.0 against 33.3\,ms, at 10\%). Qdrant's
gap of 0.034 (0.020 to 0.049) averages opposite effects: at 2\% its recall@10 was 0.024 for
uncorrelated questions, 0.312 for negatively and 0.808 for positively correlated ones, and at 30\%
0.818, 0.565 and 0.985 (build 0). This is consistent with a random 2\% of the graph being the hardest permitted subgraph to
traverse, since few of a node's neighbors would belong to it and topical groups would stay better connected; we did not measure connectivity.

\paragraph{Local or global selectivity (H2.2, falsified).}
Local selectivity predicted recall better than global selectivity in three stores and worse in five;
the pre-registered rule calls the hypothesis falsified once half of the assessable stores (eight;
OpenSearch had too few misses) go the other way. The split is consistent with the mechanism. Where recall depends on what lies near the question, in pgvector's scan in order of
distance, LanceDB's probed partitions and Elasticsearch, local selectivity won ($D$ of 0.244 to
0.538), and ACORN's query correlation correlated with recall more strongly still (Spearman 0.562 to
0.672). Where the store chooses its search path from the filter's estimated size, as Qdrant, Weaviate
and Redis document~\cite{qdrant2026search,weaviate2026vectorindex,redis2026vectors}, and in Milvus and
MongoDB, small filters are searched exactly and global selectivity won ($D$ of $-0.531$ to $-0.076$).
HCBGen's over-fetch estimate, which grows as the readable set shrinks, correlated \emph{positively}
with recall in eight of the nine stores (up to 0.551), the opposite sign to a measure of difficulty,
for the same reason. No single axis describes every store; the paper therefore reports recall against
several (released with the run).

\paragraph{Effort and the gap (H2.3, confirmed).}
At the top of each sweep the gap narrowed and remained in two of the four stores that had one, which meets the rule's bar of half. pgvector's fell from 0.424 to 0.172 (0.161 to 0.183) at \texttt{ef\_search} 1,000, and
LanceDB's moved by 0.001 at 500 probes, at its quantization floor, a negligible narrowing that still counts under the rule. Weaviate's closed, to 0.004 at
\texttt{ef} 1,024. Elasticsearch's widened, from 0.137 to 0.162 at 10,000 candidates: more candidates
helped uncorrelated questions more. pgvector reached recall@10 of 0.95 for uncorrelated questions at
\texttt{ef\_search} 400 and for negatively correlated ones at no level of its sweep; Weaviate needed one
more level for them (\texttt{ef} 256 against 128).

\paragraph{Intermediate filters harder than the smallest (H2.4, confirmed).}
Five stores lost more at an intermediate selectivity than at the smallest, as predicted; the test
does not ask whether recall recovers at larger selectivities. On clustered
layouts with negatively correlated questions and one disjunction, recall@10 at 0.06\%, 0.5\%, 2\% and
10\% was 0.927, 0.380, 0.000 and 0.000 for pgvector at its default (whose planner chose the exact
plan for most of the smallest filters, below), 1.000, 1.000, 0.460 and 0.489 for Qdrant, 0.841,
0.727, 0.641 and 0.607 for Elasticsearch, and 1.000 down to 0.822 and 0.896 at 10\% for Weaviate and
Milvus. Redis stayed at 1.000 and LanceDB near its floor (0.197 to 0.235); OpenSearch and MongoDB lost
at most 0.0012. In this regime very selective permission filters were therefore the easy case for five of the nine stores; the test does not tell us where recall is lowest.

\paragraph{pgvector's plans.}
With each search planned with its own values, PostgreSQL chose the HNSW scan for 98.6\% of the
enterprise searches at \texttt{ef\_search} 40 to 200, as in X1's custom plans, but for 88.9\% at 400
and none at 1,000, where its cost estimate made the exact plan cheaper: recall@10 1.000 at a median
of 169.7\,ms, against 0.670 at 17.9\,ms with \texttt{ef\_search} 200. On MS MARCO, at the default
setting, it chose the exact plan for 86\% to 94\% of the searches at 0.06\% selectivity, for up to
39\% at 0.5\%, and for none above.
These switches are the planner's, made per search; the plan cache of Section~\ref{sec:qual-pgvector},
which decides once per statement and connection, had no role here.

\paragraph{The enterprise corpus.}
The enterprise results agree with X1 and the audit. Qdrant reached 0.663 at \texttt{hnsw\_ef} 128,
0.864 at 1,024, 0.830 with ACORN at its default gate and 0.953 with ACORN at gate 1.0 and
\texttt{hnsw\_ef} 512, and 0.941 with per-tenant graphs at \texttt{hnsw\_ef} 128 at a similar latency
(4.7 against 4.9\,ms). Elasticsearch's default index for 768 dimensions, \texttt{bbq\_hnsw},
lost against \texttt{int8\_hnsw} at every level (0.815 against 0.888 at the default, 0.867 against
0.966 at 500 candidates). Turning off Weaviate's flat-search cutoff lowered its recall@10 from 0.923
to 0.837 here and from 0.968 to 0.812 on MS MARCO: the exact scan of filters below 40,000 objects
carries much of its recall.
